% TOP_PROBLEM: {"id": 9700013, "title": "Random Eulerian excursion dichotomy on high-dimensional tori", "category_id": 19, "category": {"id": 19, "name": "probability", "display_name": "Probability", "description": "Problems involving probability theory, stochastic processes, and random structures.", "slug": "probability", "order_index": 19, "created_at": "2026-07-31T15:26:25.671Z"}, "status": "open", "problem_number": "AMR-096-0013", "set_id": 13}
% FIRST_POSTED: 2026-10-01
% ENTRY_KIND: statement_only
% UnsolvedMath ID 9700013; source catalogue code AMR-096-0013
% !TeX program = lualatex
% Definition and sources only; this file is not an LLM solution attempt.
\documentclass[11pt,a4paper]{article}
\usepackage[margin=25mm]{geometry}
\usepackage{fontspec}
\setmainfont{lmroman10-regular.otf}[BoldFont=lmroman10-bold.otf,
 ItalicFont=lmroman10-italic.otf,BoldItalicFont=lmroman10-bolditalic.otf]
\usepackage{amsmath,amssymb,mathtools}
\usepackage{unicode-math}
\setmathfont{latinmodern-math.otf}
\AtBeginDocument{\let\setminus\smallsetminus}
\usepackage{xurl,hyperref}
\hypersetup{colorlinks=true,urlcolor=blue}
\setcounter{secnumdepth}{0}
\setlength{\parindent}{0pt}
\setlength{\parskip}{5pt}
\setlength{\emergencystretch}{3em}
\begin{document}
\section*{Random Eulerian excursion dichotomy on high-dimensional tori}\label{9700013}
{\small UnsolvedMath ID 9700013 · AMR-096-0013 · Probability · AMR Open Problem Lists}
\subsection{Definitions and mathematical statement}
Fix $d\ge3$. For each $N\ge3$, form the nearest-neighbor torus $(\mathbb Z/N\mathbb Z)^d$ and replace each undirected edge by its two opposite directed arcs. Choose uniformly an Eulerian circuit, meaning a closed traversal using every directed arc exactly once. Root its traversal at the origin; choosing uniformly among origin-rooted traversals gives the same law for the unordered excursion lengths.

Cut the circuit at successive returns to the origin. There are $2d$ excursions, and their lengths, measured in directed steps, sum to $2dN^d$. For a deterministic threshold $\omega_N\to\infty$ with $\omega_N^2<N^d$, let $b_N$ count lengths greater than $N^d/\omega_N$, let $t_N$ count lengths smaller than $\omega_N$, and let $m_N$ count the remaining lengths, including any endpoint equalities.

Is there a probability law $S^*$ on $\{1,\ldots,2d\}$ such that, for thresholds tending to infinity sufficiently slowly,
\[
(b_N,t_N,m_N)\xrightarrow{d}(S^*,2d-S^*,0)?
\]
One precise meaning of ``sufficiently slowly'' is that some deterministic $a_N\to\infty$, with $a_N=o(N^{d/2})$, works for every divergent $\omega_N\le a_N$ eventually.

Convergence in distribution here is convergence of probabilities on the finite set of integer triples. The limiting law may depend on $d$, but not on the chosen slow threshold. The target asserts that intermediate-sized excursions disappear and that a positive number remain on the macroscopic scale. It does not prescribe the exact law of the macroscopic lengths themselves.
\subsection{Short English statement}
Do random Eulerian excursions on higher-dimensional tori split into microscopic and macroscopic groups with no intermediate lengths?
\subsection{Sources}
\begin{itemize}
\item[S1] ULAM.AI and the UnsolvedMath Contributors, supplied problems.json, record 9700013 (AMR-096-0013), AMR Open Problem Lists. Catalogue identity and source transcription; CC BY 4.0.
\url{https://huggingface.co/datasets/ulamai/UnsolvedMath}.
\item[S2] Aldous - Open problems, Random Eulerian Circuits, Conjecture 0.1. Original source indexed by AMR-096-0013.
\url{https://www.stat.berkeley.edu/~aldous/Research/OP/index.html}.
\item[S3] D. Aldous and J. Yu, Random Eulerian Circuits (2014), original numbered conjecture and examples.
\url{https://www.stat.berkeley.edu/~aldous/Papers/aldous_eulerian.pdf}.
\end{itemize}
\end{document}
